\newpage
\section{The NNPDF3.1 global analysis}

\label{sec:results}

We now present the results of the NNPDF3.1 global
analysis at LO, NLO and NNLO, and compare them with the previous
release NNPDF3.0 and  with other recent PDF sets.
%
Here we present results obtained using  the complete dataset
of Tab.~\ref{tab:completedataset}-\ref{tab:completedataset3}, discussed 
in Sect.~\ref{sec:expdata}. Studies of the impact of individual
measurements will be discussed along with PDF determinations from reduced
datasets in Sect.~\ref{sec:impactnewdata}. 

After a brief methodological summary, 
we discuss the  fit quality, and 
then examine individual PDFs and their uncertainties.
We compare NNPDF3.1 PDFs with NNPDF3.0 and with  CT14~\cite{Dulat:2015mca}, MMHT2014~\cite{Harland-Lang:2014zoa} 
and ABMP16~\cite{Alekhin:2017kpj}. 
We next examine the impact of independently parametrizing
charm, the principal methodological 
improvement in NNPDF3.1. Finally, we discuss theoretical
uncertainties, both related to QCD parameters 
and to missing higher order corrections to the
theory used for PDF determination.

In this Section all NLO and NNLO NNPDF3.1 results are produced using
the CMC~\cite{Carrazza:2015hva} optimized 100 replica Monte Carlo sets,
see Sect.~\ref{sec:delivery} below: despite only
including 100 replicas, these sets reproduce  the statistical features
of a set of at least about 400 replicas (see 
Sect.~\ref{sec:compression}).  We present here only a selection of
results: a more extensive set of results is available from a public
repository, see  Sect.~\ref{sec:delivery}.

\subsection{Methodology}
\label{sec:methodology}

NNPDF3.1 PDFs are determined with largely the same methodology as in NNPDF3.0:
the only significant change is that now charm is independently
parametrized. The PDF parametrization is identical to that
discussed in Sect.~3.2 of Ref.~\cite{Ball:2014uwa}, including the
treatment of preprocessing, but with the PDF
basis in Eq.~(3.4) of that reference now supplemented by an extra PDF
for charm, parametrized like all other PDFs (as per Eq.~(2) of
Ref.~\cite{Ball:2016neh}).
PDFs are parametrized at the scale $Q_0=1.65$~GeV whenever the charm PDF
is independently parametrized. For the purposes of comparison we also provide PDF sets constructed with perturbatively 
generated charm; 
in these sets, PDFs are parametrized at the scale $Q_0=1.0$~GeV. This ensures  
that the parametrization scale is always above the charm mass when
charm is independently parametrized, and below it when it is
perturbatively generated. 

As in Ref.~\cite{Ball:2014uwa} we use $\alpha_s(m_Z)=0.118$ as a
default throughout the paper, though determinations have also been
performed for several different values of
$\alpha_s$ (see Sect.~\ref{sec:delivery}).
Pole heavy quark masses are used throughout, with the main motivation
that for the inclusive observables used for PDF determination
$\overline{\rm MS}$ masses are inappropriate, since they distort the perturbative
expansion in the threshold region~\cite{Ball:2016qeg}. 
The default values of the heavy quark pole masses are  $m_c=1.51$~GeV for
  charm and $m_b=4.92$~GeV for bottom,
  following the recommendation of the Higgs cross-section working
  group~\cite{deFlorian:2016spz};
  PDF sets for different charm mass values, corresponding to
 the $\pm 1$-sigma uncertainty band from
  Ref.~\cite{deFlorian:2016spz}, are also provided, see
Sect.~\ref{sec:delivery}.


\subsection{Fit quality}
\label{sec:fquality}

%%%%%%%%%%%%%%%%%%%%%
%%%%%%%%%%%%%%%%%%%%%%%%%%%%%%%%%%%%%%%%%%%%%%%%%%%%%%%%%%%%%%%%%%%%%%%%%%%%%%%%%%
\begin{table}[h]
\begin{center}
  \scriptsize
      \renewcommand{\arraystretch}{1.10}
\begin{tabular}{|l|c|c|c||c|c|}
\hline
  & \multicolumn{3}{c||}{NNPDF3.1  } & \multicolumn{2}{c|}{NNPDF3.0 } 
\\Dataset  &  NNLO  & NLO & LO  & NNLO  & NLO  \\
\hline
\hline
NMC    &    1.30     &   1.35   &     3.25       &   1.29   &   1.36  \\    
SLAC    &    0.75      &   1.17   &   3.35       &   0.66   &   1.08  \\    
BCDMS    &    1.21      &   1.17   &  2.20        &   1.31   &   1.21  \\    
CHORUS    &    1.11      &   1.06   &   1.16      &   1.11   &   1.14  \\    
NuTeV dimuon    &    0.82      &   0.87   &  4.75         &   0.69  &   0.61  \\    
\hline
HERA I+II inclusive    &    1.16      &   1.14   &  1.77         &   1.25   &   1.20  \\      
HERA $\sigma_c^{\rm NC}$    &    1.45      &   1.15 ()   &  1.21           &  [1.61]   &   [2.57]  \\    
HERA $F_2^b$    &    1.11      &   1.08   &   11.2           &   [1.13]   &   [1.12]  \\    
\hline
\hline
DY E866 $\sigma^d_{\rm DY}/\sigma^p_{\rm DY}$ &    0.41     &   0.40   &     1.06       &   0.47   &   0.53  \\    
DY E886 $\sigma^p$    &    1.43     &   1.05   &    0.81       &   1.69   &   1.17  \\    
DY E605  $\sigma^p$   &    1.21     &   0.97   &  0.66      &   1.09   &   0.87  \\    
\hline
CDF $Z$ rap    &    1.48     &   1.619   &    1.54         &   1.55   &   1.28  \\    
CDF Run II $k_t$ jets    &    0.87      &   0.84   &  1.07           &   0.82   &   0.95  \\    
\hline
D0 $Z$ rap    &    0.60      &   0.67   &    0.65         &   0.61   &   0.59  \\    
D0 $W\to e\nu$  asy   &    2.70      &   1.59   &   1.75         &   [2.68]   &   [4.58]  \\    
D0 $W\to \mu\nu$  asy    &    1.56      &   1.52   &    2.16          &   [2.02]   &   [1.43]  \\    
\hline
\hline
ATLAS total   &    {\bf 1.09}      &   {\bf 1.36}   &    {\bf 5.34}         &   {\bf 1.92 }  & {\bf 1.98}  \\    
ATLAS $W,Z$ 7 TeV 2010    &    0.96     &   1.04   &    2.38         &   1.42   &   1.39  \\    
ATLAS high-mass DY 7 TeV    &    1.54     &   1.88   &   4.05         &   1.60   &   2.17  \\    
ATLAS low-mass DY 2011    &    0.90      &   0.69   &   2.86         &   [0.94]   &   [0.81]  \\    
ATLAS $W,Z$ 7 TeV 2011    &    2.14     &   3.70   &    27.2       &   [8.44]   &   [7.6]  \\    
ATLAS jets 2010 7 TeV     &    0.94      &   0.92   &  1.22        &   1.12   &   1.07  \\    
ATLAS jets 2.76 TeV     &    1.03      &   1.03   &     1.50     &   1.31   &   1.32  \\    
ATLAS jets 2011 7 TeV     &    1.07     &   1.12   &   1.59       &   [1.03]   &   [1.12]  \\    
ATLAS $Z$ $p_T$ 8 TeV $(p_T^{ll},M_{ll})$       &    0.93    &   1.17   &   -         &   [1.05]  &  [1.28]  \\    
ATLAS $Z$ $p_T$ 8 TeV $(p_T^{ll},y_{ll})$    &    0.94      &   1.77   &   -         &   [1.19]   &   [2.49]  \\    
ATLAS $\sigma_{tt}^{\rm tot}$     &    0.86      &   1.92   &  53.2           &   0.67   &   1.07  \\    
ATLAS $t\bar{t}$ rap    &    1.45     &   1.31   &  1.99       &   [3.32]   &   [1.50]  \\
\hline
CMS total   &   {\bf  1.06}    &   {\bf 1.20 }  &    {\bf 2.13}        &   {\bf 1.19 }  &
{\bf 1.33}  \\    
CMS $W$ asy 840 pb     &    0.78     &   0.86   &    1.55       &   0.73   &   0.85  \\    
CMS $W$ asy 4.7 fb     &    1.75     &   1.77   &   3.16      &   1.75   &   1.82  \\    
CMS $W+c$ tot    &    -     &   0.54   &    16.5          &   -   &   0.93  \\    
CMS $W+c$ ratio    &    -      &   1.91   &    3.21         &   -   &   2.09  \\    
CMS Drell-Yan 2D 2011    &    1.27     &   1.23   &  2.15           &   1.20   &   1.19  \\    
CMS $W$ rap 8 TeV    &    1.01      &   0.70   &    4.32         &   [1.24]   &   [0.96]  \\    
CMS jets 7 TeV 2011     &    0.84      &   0.84   &  0.93            &   1.06   &   0.98  \\    
CMS jets 2.76 TeV    &    1.03      &   1.01   &    1.09          &   [1.22]   &   [1.18]  \\    
CMS $Z$ $p_T$ 8 TeV $(p_T^{ll},M_{ll})$    &    1.32    &   3.65   &   -         &   [1.59]   &   [3.86]  \\    
CMS $\sigma_{tt}^{\rm tot}$     &    0.20      &   0.59   &  53.4            &   0.56   &   0.10  \\    
CMS $t\bar{t}$ rap    &    0.94      &   0.96   &     1.32           &   [1.15]   &   [1.01]  \\
\hline 
LHCb total   &  {\bf  1.47}      &   {\bf 1.62}   &   {\bf 5.16}       &   {\bf 2.11}   &   {\bf 2.67}  \\
LHCb $Z$ 940 pb    &    1.49      &   1.27   &        2.51      &   1.29   &   0.91  \\    
LHCb $Z\to ee$ 2 fb    &    1.14      &   1.33   &    6.34         &   1.21   &   2.31  \\    
LHCb $W,Z \to \mu$ 7 TeV    &    1.76      &   1.60   &   4.70          &   [2.59]   &   [2.36]  \\    
LHCb $W,Z \to \mu$ 8 TeV    &    1.37     &   1.88   &   7.41            &   [2.40]   &   [3.74]  \\    
\hline
\hline
{\bf Total dataset}   &  {\bf  1.148}      & {\bf  1.168}   & {\bf  2.238}        & {\bf  1.284}   & {\bf  1.307}  \\    
\hline
\end{tabular}
\caption{\small
\label{tab:chi2tab_31-nlo-nnlo-30}
  The values of $\chi^2/N_{\rm dat}$ for the global fit and for 
all the datasets included in the NNPDF3.1 LO, NLO
  and NNLO PDF determinations. Values obtained using the NNPDF3.0 NLO and NNLO
  PDFs  are also shown: numbers in brackets correspond 
   to data not fitted in NNPDF3.0. Note that NNPDF3.0 values are
   produced using NNPDF3.1 theory settings, and are thus somewhat
   worse than those quoted in Ref.~\cite{Ball:2014uwa}.
}
\end{center}
\end{table}
%%%%%%%%%%%%%%%%%%%%%%%%%%%

%%%%%%%%%%%%%%%%%%%%%


In Table~\ref{tab:chi2tab_31-nlo-nnlo-30} we provide
values of $\chi^2/N_{\rm dat}$ both for the global fit and individually for all the datasets included in the
NNPDF3.1 LO, NLO
  and NNLO PDF determinations. These are compared to their NNPDF3.0 NLO and NNLO
  counterparts. The $\chi^2$
  is computed using the covariance matrix including all correlations,
  as published by the corresponding experiments.
 Inspection of this table
shows that the fit quality improves from LO to NLO to NNLO: not only
is there a significant improvement between LO and NLO, but there is also a marked improvement when going from NLO
to NNLO. It is interesting to note that this was not the case in
NNPDF3.0 where the fit quality at NNLO was in fact slightly worse than
at NLO (see Table~9 of Ref.~\cite{Ball:2014uwa}). This reflects the increased proportion of 
hadronic processes included in NNPDF3.1, for which NNLO corrections
are often substantial, and also, possibly, methodological improvements.

The overall fit quality with NNPDF3.1 is  rather better than that
obtained using NNPDF3.0 PDFs. 
Whereas this is clearly expected for LHC measurements which
were not included in NNPDF3.0, it is
interesting to note that the HERA
measurements which were already present in 3.0 (though in slightly different
uncombined form) are also better fitted. 
 The quality of the description with the
previous NNPDF3.0 PDFs is nevertheless quite acceptable for all the new data, 
indicating a general compatibility between NNPDF3.0 and NNPDF3.1. 
Note that NNPDF3.0 values in
Table~\ref{tab:chi2tab_31-nlo-nnlo-30} are computed using the NNPDF3.1
theory settings, thus in particular with different 
values of the heavy quark masses than those used in the NNPDF3.0 PDF
determination. Because of this, the NNPDF3.0 fit quality
shown shown in Table~9 of Ref.~\cite{Ball:2014uwa} is slightly better than 
that shown in Table~\ref{tab:chi2tab_31-nlo-nnlo-30}, yet even so 
the fit quality of NNPDF3.1 is better still. Specifically, 
concerning HERA data, the fit quality of
NNPDF3.0 with consistent theory settings can be read off Table~7 of
Ref.~\cite{Czakon:2016olj}: it corresponds to  $\chi^2/N_{\rm
  dat}=1.21$ thereby showing that indeed NNPDF3.1 provides a better
description. The reasons for this improvement will be discussed in
Sect.~\ref{sec:results-mc} below.

For many of the new LHC measurements, achieving a good description of the data
is only possible at NNLO. The
total $\chi^2/N_{\rm dat}$ for the ATLAS, CMS and
LHCb experiments is 1.09, 1.06 and 1.47 respectively at NNLO,
compared to 1.36, 1.20 and 1.62 at NLO. The datasets exhibiting
the largest improvement when going from NLO to NNLO are those with the
smallest experimental uncertainties. For example the
ATLAS $W,Z$ 2011 rapidity distributions
(from 3.70 to 2.14), the CMS 8 TeV $Z$ $p_T$ distributions
(from 3.65 to 1.32) and the LHCb 8 TeV $W,Z\to \mu$
rapidity distributions (from 1.88 to 1.37); in these experiments
uncorrelated statistical uncertainties are typically at the
sub-percent level.
%
It is likely that this trend will continue as LHC measurements become more precise.


% %%%%%%%%%%%%%%%%%%%%%
% \begin{table}[t]
\begin{center}
  \scriptsize
  \renewcommand{\arraystretch}{1.10}
\begin{tabular}{|l|c|c||c|c|}
\hline
 & \multicolumn{2}{c|}{NNPDF3.1 NNLO } & \multicolumn{2}{c|}{NNPDF3.1 NLO } \\
Dataset & $\chi^2_{\rm exp}$ & $\chi^2_{t_0}$ & $\chi^2_{\rm exp}$  & $\chi^2_{t_0}$    \\
\hline
\hline
NMC    &    1.30   &   1.23   &   1.35   &   1.35  \\    
SLAC    &    0.75   &   0.73   &   1.17   &   1.12  \\    
BCDMS    &    1.21   &   1.23   &   1.17   &   1.20  \\    
CHORUS    &    1.11   &   1.14   &   1.06   &   1.09  \\    
NuTeV dimuon    &    0.82   &   0.76   &   0.87   &   0.81  \\
\hline
HERA I+II inclusive    &    1.16   &   1.26   &   1.14   &   1.29  \\    
HERA $\sigma_c^{\rm NC}$    &    1.45   &   1.46   &   1.15   &   1.15  \\    
HERA $F_2^b$   &    1.11   &   1.33   &   1.08   &   1.27  \\    
\hline
\hline
DYE866 $\sigma^d_{\rm DY}/\sigma^p_{\rm DY}$    &    0.41   &   0.41   &   0.40   &   0.40  \\    
DYE866 $\sigma^p$     &    1.43   &   1.09   &   1.05   &   0.98  \\    
DYE605    &    1.21   &   1.02   &   0.99   &   1.04  \\    
\hline
CDF $Z$ rap     &    1.48   &   1.46   &   1.62   &   1.76  \\    
CDF Run II $k_t$ jets    &    0.87   &   1.03   &   0.84   &   1.02  \\
\hline
D0 $Z$ rap    &    0.60   &   0.61   &   0.66   &   0.68  \\    
D0 $W\to e\nu$  asy    &    2.71   &   2.71   &   1.59   &   1.59  \\    
D0 $W\to \mu\nu$  asy    &    1.56   &   1.56   &   1.52   &   1.52  \\
\hline
\hline
ATLAS total   &    1.09   &   0.99   &   1.36   &   1.37  \\    
ATLAS $W,Z$ 7 TeV 2010    &    0.96   &   0.88   &   1.04   &   1.00  \\    
ATLAS HM DY 7 TeV    &    1.54   &   1.76   &   1.88   &   2.33  \\    
ATLAS low-mass DY 2011    &    0.90   &   0.99   &   0.69   &   0.78  \\    
ATLAS $W,Z$ 7 TeV 2011    &    2.14   &   2.19   &   3.70   &   4.02  \\    
ATLAS jets 2010 7 TeV     &    0.95   &   0.51   &   0.92   &   0.51  \\    
ATLAS jets 2.76 TeV     &    1.03   &   0.58   &   1.03   &   0.60  \\    
ATLAS jets 2011 7 TeV     &    1.07   &   1.51   &   1.12   &   1.57  \\    
ATLAS $Z$ $p_T$ 8 TeV $(y_{ll},M_{ll})$       &    0.93   &   1.04   &   1.17   &   1.57  \\    
ATLAS $Z$ $p_T$ 8 TeV $(p_T^{ll},M_{ll})$    &    0.94   &   1.13   &   1.76   &   2.35  \\    
ATLAS $\sigma_{tt}^{\rm tot}$     &    0.86   &   0.93   &   1.92   &   2.19  \\    
ATLAS $t\bar{t}$ rap    &    1.45   &   1.47   &   1.31   &   1.33  \\
\hline
CMS total   &    1.06   &   1.19   &   1.20   &   1.43  \\    
CMS $W$ asy 840 pb     &    0.78   &   0.78   &   0.86   &   0.86  \\    
CMS $W$ asy 4.7 fb     &    1.75   &   1.75   &   1.77   &   1.77  \\    
CMS $W+c$ tot    &    -   &   -   &   0.54   &   0.65  \\    
CMS $W+c$ ratio    &    -   &   -   &   1.91   &   1.91  \\    
CMS Drell-Yan 2D 2011    &    1.27   &   1.44   &   1.23   &   1.43  \\    
CMS $W$ rap 8 TeV    &    1.01   &   0.95   &   0.70   &   0.74  \\    
CMS jets 7 TeV 2011     &    0.84   &   0.90   &   0.84   &   0.89  \\    
CMS jets 2.76 TeV    &    1.03   &   1.31   &   1.01   &   1.28  \\    
CMS $Z$ $p_T$ 8 TeV $(p_T^{ll},M_{ll})$    &    1.32   &   1.46   &   3.65   &   5.13  \\    
CMS $\sigma_{tt}^{\rm tot}$     &    0.20   &   0.19  &   0.59   &   0.63  \\    
CMS $t\bar{t}$ rap    &    0.94   &   0.96   &   0.96   &   0.96  \\
\hline
LHCb total   &    1.47   &   1.50   &   1.61   &   1.80  \\    
LHCb $Z$ 940 pb     &    1.49   &   1.67   &   1.27   &   1.48  \\    
LHCb $Z\to ee$ 2 fb     &    1.14   &   1.16   &   1.33   &   1.44  \\    
LHCb $W,Z \to \mu$ 7 TeV    &    1.76   &   1.80   &   1.60   &   1.75  \\    
LHCb $W,Z \to \mu$ 8 TeV    &    1.37   &   1.35   &   1.88   &   2.14  \\    
\hline
\hline
Total dataset     &  {\bf  1.148}   &{\bf   1.183}   & {\bf  1.168}   & {\bf  1.246}  \\    
\hline
\end{tabular}
\caption{\small
  \label{tab:chi2tab_exp_vs_t0}
  Same as Table~\ref{tab:chi2tab_31-nlo-nnlo-30} now comparing the values of
  the  $\chi^2/N_{\rm dat}$ obtained with either  
  the experimental or the $t_0$ definition (see text), for the NNPDF3.1 NLO and NNLO fits.
}
\end{center}
\end{table}
%%%%%%%%%%%%%%%%%%%%%%%%%%%%%%%%%%%%%%%%%%%%%%%%%%%%%%%%%%%%%%%
%%%%%%%%%%%%%%%%%%%%%%%%%%%%%%%%%%%%%%%%%%%%%%%%%%%%%%%%%%%%%%%


% %%%%%%%%%%%%%%%%%%%%%

% In Table~\ref{tab:chi2tab_exp_vs_t0} we compare
%  $\chi^2/N_{\rm dat}$ values obtained using two different definitions
% of this figure of merit; namely using the covariance matrix as
% published by the respective experiments (``experimental'' $\chi^2$)
% and an alternative definition in which the covariance matrix is
% computed using the theoretical prediction obtained using the results
% of a previous fit (``$t_0$'' $\chi^2$). The latter definition 
% (see Refs.~\cite{Ball:2009qv,Ball:2012wy}) is
% the one used for minimization, as the experimental $\chi^2$ would lead to biased
% results in the presence of multiplicative
% uncertainties~\cite{D'Agostini:1993uj}. It is clear from the table
% that $\chi^2$ values determined using either of these two definitions
% are reassuringly similar, and that in many cases $\chi^2_{\rm exp}$ turns 
% out to be lower than the $\chi^2_{t_0}$.

\clearpage

\subsection{Parton distributions}
\label{sec:PDFcomparisons}

We now inspect the baseline NNPDF3.1 parton distributions, and
compare them to NNPDF3.0 
and to  MMHT14~\cite{Harland-Lang:2014zoa},
CT14~\cite{Dulat:2015mca} and
ABMP16~\cite{Alekhin:2017kpj}.
%
The NNLO NNPDF3.1 PDFs are displayed in Fig.~\ref{fig:nnlopdfs}. It can be seen 
that although charm is now independently parametrized, it is still known more precisely than the strange PDF.
The most precisely determined PDF over most of the experimentally accessible range of $x$ is now 
the gluon, as will be discussed in more detail below. 


%%%%%%%%%%%%%%%%%%%%%%%%%%%%%%%%%%%%%%%%%%%%%%%%%%%%%%%%%%%%%%%%%%%%%
\begin{figure}[t]
\begin{center}
  \includegraphics[scale=0.75]{images/plots/nnpdf31nnlo-10.pdf}
  \includegraphics[scale=0.75]{images/plots/nnpdf31nnlo-1e4.pdf}
  \caption{\small The NNPDF3.1 NNLO PDFs, evaluated
    at $\mu^2=10~{\rm GeV}^2$ (left) and $\mu^2=10^4~{\rm GeV}^2$ (right). 
    \label{fig:nnlopdfs}
  }
\end{center}
\end{figure}
%%%%%%%%%%%%%%%%%%%%%%%%%%%%%%%%%%%%%%%%%%%%%%%%%%%%%%%%%%%%%%%%%%%%%%

 
In Fig.~\ref{fig:distances_31_vs_30} we show the  distance between the NNPDF3.1 and NNPDF3.0 PDFs.
According to  the definition of the distance given in Ref.~\cite{Ball:2010de}, $d\simeq 1$ corresponds to
statistically equivalent sets. Comparing two sets with $N_{\rm rep}=100$ replicas,
a distance of $d\simeq 10$ corresponds to a difference of one-sigma in units of the corresponding variance,
both for central values and for PDF uncertainties. For clarity only the distance between the total
strangeness distributions $s^+=s+\bar s$ is shown, rather than the strange and
antistrange separately.
%
We find important differences both at the level of central values and
of PDF errors for all flavors and in the entire range of $x$.
%
The largest distance is found for charm, which is independently
parametrized in NNPDF3.1,
while it was not in NNPDF3.0. Aside from this, the most significant
distances are seen in 
light quark distributions at large $x$ and strangeness at medium $x$. 

%%%%%%%%%%%%%%%%%%%%%%%%%%%%%%%%%%%%%%%%%%%%%%%%%%%%%%%%%%%%%%%%%%%%%
\begin{figure}[t]
\begin{center}
  \includegraphics[scale=1]{images/plots/distances_31_vs_30.pdf}
  \caption{\small Distances between the central
    values (left) and the uncertainties (right) of
    the NNPDF3.0 and NNPDF3.1 NNLO PDF sets, evaluated
    at $Q=100$ GeV. Note the different in scale on the $y$ axis between
    the two plots.
    \label{fig:distances_31_vs_30}
  }
\end{center}
\end{figure}
%%%%%%%%%%%%%%%%%%%%%%%%%%%%%%%%%%%%%%%%%%%%%%%%%%%%%%%%%%%%%%%%%%%%%%

In Fig.~\ref{fig:31-nnlo-vs30} we compare the full set of NNPDF3.1 NNLO PDFs with NNPDF3.0.
%
The NNPDF3.1 gluon is slightly larger than its NNPDF3.0 counterpart
in the $x\lsim 0.03$ region, while it becomes smaller at larger $x$,
with significantly reduced PDF errors.
%
The NNPDF3.1 light quarks and strangeness are larger than 3.0
at intermediate $x$, with the largest deviation seen for the strange
and antidown PDFs, while at both small and large $x$
there is good agreement between the two PDF determinations. The best-fit
charm PDF of NNPDF3.1 is significantly smaller in the intermediate-$x$ region 
compared to the perturbative charm of NNPDF3.0, while at larger $x$ it has 
significantly increased uncertainty.

A detailed comparison of the corresponding uncertainties is presented
in Fig.~\ref{fig:ERR-31-nnlo-vs30}, where we compare the relative
uncertainty on each PDF, defined as the ratio of
the one-sigma PDF uncertainty to the central value of the
NNPDF3.1 set. NNPDF3.1 uncertainties
are either comparable to those of NNPDF3.0, or are rather smaller. The only 
major exception to this is the charm PDF at intermediate and large $x$ for 
which uncertainties are substantially increased. On the other hand, the uncertainties 
in the gluon PDF are smaller in NNPDF3.1 over the entire range of $x$.
This is an important result, since one may have expected generally larger
uncertainties in NNPDF3.1 due to the inclusion of one additional
freely parametrized PDF. 
The fact that the only uncertainty which has enlarged significantly 
is that of the charm PDF suggests that not parametrizing charm may be a source of bias. 
The fact that central values change by a non-negligible amount, though
 compatible within uncertainties, while the uncertainties themselves are significantly
reduced, strongly suggests that NNPDF3.1 is more accurate than
 NNPDF3.0, as would be expected from the substantial amount of new data 
 included in the fit.
The effect of parametrizing charm on PDFs and their uncertainties will be
discussed in more detail in Sect.~\ref{sec:results-mc}, while the effects of 
the new data on both central values and uncertainties will be
discussed in Sect.~\ref{sec:disentangling}.

%%%%%%%%%%%%%%%%%%%%%%%%%%%%%%%%%%%%%%%%%%%%%%%%%%%%%%%%%%%%%%%%%%%%%
\begin{figure}[t]
  \begin{center}
      \includegraphics[scale=0.32]{images/plots/xu-31-nnlo-vs30.pdf}
    \includegraphics[scale=0.32]{images/plots/xubar-31-nnlo-vs30.pdf}
    \includegraphics[scale=0.32]{images/plots/xd-31-nnlo-vs30.pdf}
    \includegraphics[scale=0.32]{images/plots/xdbar-31-nnlo-vs30.pdf}
  \includegraphics[scale=0.32]{images/plots/xs-31-nnlo-vs30.pdf}
  \includegraphics[scale=0.32]{images/plots/xsbar-31-nnlo-vs30.pdf}
  \includegraphics[scale=0.32]{images/plots/xc-31-nnlo-vs30.pdf}
  \includegraphics[scale=0.32]{images/plots/xg-31-nnlo-vs30.pdf}
    \caption{\small Comparison between NNPDF3.1 and NNPDF3.0 NNLO PDFs
      at $Q=100$ GeV. From top to bottom  up and antiup, down
      and antidown, strange and antistrange, charm and gluon are shown.
    \label{fig:31-nnlo-vs30}
  }
\end{center}
\end{figure}
%%%%%%%%%%%%%%%%%%%%%%%%%%%%%%%%%%%%%%%%%%%%%%%%%%%%%%%%%%%%%%%%%%%%%%




%%%%%%%%%%%%%%%%%%%%%%%%%%%%%%%%%%%%%%%%%%%%%%%%%%%%%%%%%%%%%%%%%%%%%
\begin{figure}[t]
  \begin{center}
    \includegraphics[scale=0.32]{images/plots/xu-ERR-31-nnlo-vs30.pdf}
    \includegraphics[scale=0.32]{images/plots/xubar-ERR-31-nnlo-vs30.pdf}
    \includegraphics[scale=0.32]{images/plots/xd-ERR-31-nnlo-vs30.pdf}
    \includegraphics[scale=0.32]{images/plots/xdbar-ERR-31-nnlo-vs30.pdf}
  \includegraphics[scale=0.32]{images/plots/xs-ERR-31-nnlo-vs30.pdf}
  \includegraphics[scale=0.32]{images/plots/xsbar-ERR-31-nnlo-vs30.pdf}
  \includegraphics[scale=0.32]{images/plots/xc-ERR-31-nnlo-vs30.pdf}
  \includegraphics[scale=0.32]{images/plots/xg-ERR-31-nnlo-vs30.pdf}
  \caption{\small Comparison between NNPDF3.1 and NNPDF3.0 relative
    PDF uncertainties at $Q=100$; the PDFs are as in
    Fig.~\ref{fig:31-nnlo-vs30}. The uncertainties shown are all
    normalized to the NNPDF3.1 central value.
    \label{fig:ERR-31-nnlo-vs30}
  }
\end{center}
\end{figure}
%%%%%%%%%%%%%%%%%%%%%%%%%%%%%%%%%%%%%%%%%%%%%%%%%%%%%%%%%%%%%%%%%%%%%%


In Fig.~\ref{fig:globalfits} we compare the NNPDF3.1 PDFs to the other
global PDF sets included in the PDF4LHC15 combination along with NNPDF3.0, 
namely CT14 and MMHT14. This comparison is therefore indicative of the effect of
replacing NNPDF3.0 with NNPDF3.1 in the combination. 
The relative 
uncertainties in the three sets are compared in
Fig.~\ref{fig:ERR-globalfits}. Comparing Fig.~\ref{fig:globalfits} to
Fig.~\ref{fig:31-nnlo-vs30},  it is interesting to observe that  
several aspects of the pattern
of differences between NNPDF3.1 and the other global fits
are  similar to those between NNPDF3.1  and NNPDF3.0, and therefore
they are likely to have a similar origin. This is patricularly clear
for the charm and gluon.
The gluon in the region $x\lsim 0.03$, relevant for Higgs
production, is still in good 
    agreement between the three sets. However, now NNPDF3.1 is at the upper edge
    of the one-sigma range, i.e. the  NNPDF3.1 gluon
    in this region is enhanced.
    %
    At large $x$ the NNPDF3.1 gluon is instead suppressed in comparison to
    MMHT14 and CT14. As we will show in Sects.~\ref{sec:results-mc}, \ref{sec:disentangling}
    the  enhancement is a consequence of parametrizing charm, while
    as we will show in 
    Sect.~\ref{sec:impacttop} the large-$x$ suppression is a direct
    consequence of including the 8~TeV top differential data. 
    The uncertainty in the NNPDF3.1 gluon PDF is now 
   noticeably smaller than that of either CT14 or MMHT14. 

For the quark PDFs, for up and down we find 
good agreement in the entire range of $x$. For the antidown PDF,
agreement is marginal, 
with NNPDF3.1 above MMHT14 and CT14 for $x\lsim 0.1$ and below them for larger
$x$. The strange fraction of the proton is larger in NNPDF3.1 than CT14 and 
MMHT14, and has rather smaller PDF uncertainties.
The best-fit NNPDF3.1 charm is suppressed at intermediate $x$ in comparison to the
perturbatively generated ones of CT14 and MMHT14, but has a much larger
uncertainty at large $x$ as would be expected, with the differences
clearly traceable  to the fact that in
   NNPDF3.1 charm  is freely parametrized.

%%%%%%%%%%%%%%%%%%%%%%%%%%%%%%%%%%%%%%%%%%%%%%%%%%%%%%%%%%%%%%%%%%%%%
\begin{figure}[t]
  \begin{center}
 \includegraphics[scale=0.32]{images/plots/xu-31-nnlo-globalfits.pdf}
    \includegraphics[scale=0.32]{images/plots/xubar-31-nnlo-globalfits.pdf}
    \includegraphics[scale=0.32]{images/plots/xd-31-nnlo-globalfits.pdf}
    \includegraphics[scale=0.32]{images/plots/xdbar-31-nnlo-globalfits.pdf}
  \includegraphics[scale=0.32]{images/plots/xs-31-nnlo-globalfits.pdf}
  \includegraphics[scale=0.32]{images/plots/xsbar-31-nnlo-globalfits.pdf}
  \includegraphics[scale=0.32]{images/plots/xc-31-nnlo-globalfits.pdf}
  \includegraphics[scale=0.32]{images/plots/xg-31-nnlo-globalfits.pdf}
  \caption{\small Comparison between NNPDF3.1, CT14 and MMHT2014 NNLO PDFs.
    The comparison is performed at $Q=100$ GeV, and results are shown
    normalized to the central value of NNPDF3.1; the PDFs are as in Fig.~\ref{fig:31-nnlo-vs30}.
    \label{fig:globalfits}
  }
\end{center}
\end{figure}
%%%%%%%%%%%%%%%%%%%%%%%%%%%%%%%%%%%%%%%%%%%%%%%%%%%%%%%%%%%%%%%%%%%%%%

%%%%%%%%%%%%%%%%%%%%%%%%%%%%%%%%%%%%%%%%%%%%%%%%%%%%%%%%%%%%%%%%%%%%%
\begin{figure}[t]
  \begin{center}
 \includegraphics[scale=0.32]{images/plots/xu-ERR-31-nnlo-globalfits.pdf}
    \includegraphics[scale=0.32]{images/plots/xubar-ERR-31-nnlo-globalfits.pdf}
    \includegraphics[scale=0.32]{images/plots/xd-ERR-31-nnlo-globalfits.pdf}
    \includegraphics[scale=0.32]{images/plots/xdbar-ERR-31-nnlo-globalfits.pdf}
  \includegraphics[scale=0.32]{images/plots/xs-ERR-31-nnlo-globalfits.pdf}
  \includegraphics[scale=0.32]{images/plots/xsbar-ERR-31-nnlo-globalfits.pdf}
  \includegraphics[scale=0.32]{images/plots/xc-ERR-31-nnlo-globalfits.pdf}
  \includegraphics[scale=0.32]{images/plots/xg-ERR-31-nnlo-globalfits.pdf}
  \caption{\small Comparison between NNPDF3.1, CT14 and MMHT2014 relative
    PDF uncertainties at $Q=100$; the PDFs are as in Fig.~\ref{fig:globalfits}.
    \label{fig:ERR-globalfits}
  }
\end{center}
\end{figure}
%%%%%%%%%%%%%%%%%%%%%%%%%%%%%%%%%%%%%%%%%%%%%%%%%%%%%%%%%%%%%%%%%%%%%%



%%%%%%%%%%%%%%%%%%%%%%%%%%%%%%%%%%%%%%%%%%%%%%%%%%%%%%%%%%%%%%%%%%%%%
\begin{figure}[t]
  \begin{center}
    \includegraphics[scale=0.32]{images/plots/xu-abmp16.pdf}
    \includegraphics[scale=0.32]{images/plots/xubar-abmp16.pdf}
    \includegraphics[scale=0.32]{images/plots/xd-abmp16.pdf}
    \includegraphics[scale=0.32]{images/plots/xdbar-abmp16.pdf}
  \includegraphics[scale=0.32]{images/plots/xs-abmp16.pdf}
  \includegraphics[scale=0.32]{images/plots/xsbar-abmp16.pdf}
  \includegraphics[scale=0.32]{images/plots/xc-abmp16.pdf}
  \includegraphics[scale=0.32]{images/plots/xg-abmp16.pdf}
  \caption{\small Same as Fig.~\ref{fig:globalfits} but now comparing
    to the ABMP16 NNLO $n_f=5$ sets both with their default
    $\alpha_s(m_Z)=0.1147$, and $\alpha_s(m_Z)=0.118$.
    \label{fig:abmp16}
  }
\end{center}
\end{figure}
%%%%%%%%%%%%%%%%%%%%%%%%%%%%%%%%%%%%%%%%%%%%%%%%%%%%%%%%%%%%%%%%%%%%%%

Finally, in Fig.~\ref{fig:abmp16} we compare NNPDF3.1 
to the recent ABMP16 set. This set is released in
various fixed-flavor number schemes. Because we perform the comparison
at a scale $Q^2=10^4$ GeV$^2$, we choose the $n_f=5$ NNLO ABMP16 sets,
both with their default value $\alpha_s(m_Z)=0.1147$ and with
$\alpha_s(m_Z)=0.118$. 
When a common value of $\alpha_s(m_Z)=0.118$ is adopted, there
is generally reasonable agreement for the gluon PDF, except
at large $x$ where ABMP16 undershoots
NNPDF3.1.
%
Differences are larger in the case of light quarks:
the ABMP16 up distribution overshoots NNPDF3.1 at large $x$, while the down quark undershoots
in the whole $x$ range. Differences are largest for the strange PDF,
though comparing with Fig.~\ref{fig:globalfits} it is clear that ABMP16
differs by a similary large amount from MMHT14 and CT14.
 In general the ABPM16
sets have rather smaller uncertainties than NNPDF3.1. This
is especially striking for strangeness, where the difference in
uncertainty is particularly evident. This is to be contrasted
with the CT14 and MMHT14 sets, which have qualitatively similar uncertainties 
to NNPDF3.1 throughout the data region. 
The fact that the uncertainties for ABMP16 are so small can be traced
to their overly restrictive parametrization, and the fact that this set is produced using a
Hessian methodology, but unlike MMHT14 and CT14, with no tolerance
(see Refs.~\cite{Harland-Lang:2014zoa,Dulat:2015mca}).



\clearpage

\clearpage

% Stability with respect charm mass
\subsection{Methodological improvements: parametrizing charm}
\label{sec:results-mc}

The main methodological improvement in NNPDF3.1 over
NNPDF3.0 is the fact that the charm PDF is now parametrized  
in the same way as the light and strange quark PDFs.
To quantify the effect of this change, we have performed a repeat of the NNPDF3.1 analysis
but with charm treated as in all previous NNPDF PDF determinations, 
i.e., generated
entirely perturbatively through matching conditions implemented at NLO or NNLO. 
%%%%%%%%%%%%%%%%%%%%
%%%%%%%%%%%%%%%%%%%%%%%%%%%%%%%%%%%%%%%%%%%%%%%%%%%%%%%%%%%%%%%%%%%%%%%%%%%%%%%%%%
\begin{table}[h]
\begin{center}
  \scriptsize
      \renewcommand{\arraystretch}{1.10}
\begin{tabular}{|l|c|c||c|c|}
  \hline
  &  \multicolumn{2}{c||}{NNPDF3.1 pert. charm} &
  \multicolumn{2}{c|}{NNPDF3.1}\\
  \hline
Dataset  &    NNLO  & NLO 
& NNLO    & NLO  \\
\hline
\hline
NMC      &   1.38     &   1.38   &    1.30     &   1.35 \\    
SLAC      &   0.70      &   1.22  &    0.75      &   1.17   \\    
BCDMS       &   1.27      &   1.24  &    1.21      &   1.17   \\    
CHORUS       &   1.10      &   1.07    &    1.11      &   1.06  \\    
NuTeV dimuon       &   1.27      &   1.01  &    0.82      &   0.87  \\
\hline
HERA I+II inclusive       &   1.21      &   1.15&    1.16      &   1.14   \\
HERA $\sigma_c^{\rm NC}$       &   1.20 (1.42)    &   1.21 (1.35)
 &    1.45      &   1.15 (1.35)\\
HERA $F_2^b$       &   1.16      &   1.12
&    1.11      &   1.08  \\
\hline
\hline
DYE866 $\sigma^d_{\rm DY}/\sigma^p_{\rm DY}$
   &   0.46     &   0.48   &    0.41     &   0.40 \\
DYE886 $\sigma^p$       &   1.38     &   1.09 &    1.43     &   1.05    \\
DYE605  $\sigma^p$      &   1.05     &   0.83 &    1.21     &   0.97  \\
\hline
CDF $Z$ rap       &   1.44      &   1.46   &    1.48     &   1.62 \\
CDF Run II $k_t$ jets       &   0.86      &   0.86
&    0.87      &   0.84 \\
\hline
D0 $Z$ rap       &   0.60      &   0.64 &    0.60      &   0.67   \\
D0 $W\to e\nu$  asy      &   2.71      &   1.63
&    2.70      &   1.59    \\
D0 $W\to \mu\nu$  asy       &   1.42      &   1.38
&    1.56      &   1.52 \\
\hline
\hline
ATLAS total      &   {\bf 1.17}     &  {\bf 1.45}  &
{\bf 1.09}      &   {\bf 1.3}  \\    
  ATLAS $W,Z$ 7 TeV 2010      &   1.04      &   1.08 &   0.96     &   1.04  \\
  ATLAS high-mass DY 7 TeV      &   1.66     &   2.08
   &    1.54     &   1.88  \\
  ATLAS low-mass DY 2011      &   0.83      &   0.70
  &    0.90      &   0.69 \\
  ATLAS $W,Z$ 7 TeV 2011       &   2.74      &   4.29
   &    2.14     &   3.70  \\
ATLAS jets 2010 7 TeV       &   0.96      &   0.95  &    0.94      &   0.92  \\
ATLAS jets 2.76 TeV       &   1.06      &   1.13   &    1.03      &   1.03  \\
ATLAS jets 2011 7 TeV       &   1.11     &   1.14   &    1.07     &   1.12  \\
ATLAS $Z$ $p_T$ 8 TeV $(p_T^{ll},M_{ll})$          &   0.94     &   1.19
&    0.93    &   1.17 \\
ATLAS $Z$ $p_T$ 8 TeV $(p_T^{ll},y_{ll})$      &   0.96      &   1.84
 &    0.94      &   1.77 \\
ATLAS $\sigma_{tt}^{\rm tot}$       &   0.80      &   2.03  &    0.86      &   1.92  \\
ATLAS $t\bar{t}$ rap      &   1.39     &   1.18  &    1.45     &   1.31  \\
\hline
CMS total     &   {\bf 1.09}     &   {\bf 1.2}
&    {\bf  1.06}    &   {\bf 1.20 } \\
CMS $W$ asy 840 pb       &   0.69     &   0.80  &    0.78     &   0.86  \\
CMS $W$ asy 4.7 fb        &   1.75      &   1.76
 &    1.75     &   1.77\\
CMS $W+c$ tot      &   -     &   0.49  &    -     &   0.54 \\
CMS $W+c$ ratio       &   -     &   1.92   &    -      &   1.91 \\
CMS Drell-Yan 2D 2011       &   1.33      &   1.27
 &    1.27     &   1.23  \\
CMS $W$ rap 8 TeV       &   0.90      &   0.65 &    1.01      &   0.70    \\
CMS jets 7 TeV 2011       &   0.87      &   0.86
&    0.84      &   0.84  \\
CMS jets 2.76 TeV       &   1.06      &   1.05  &    1.03      &   1.01  \\
CMS $Z$ $p_T$ 8 TeV $(p_T^{ll},y_{ll})$      &   1.29      &   3.50
 &    1.32    &   3.65 \\
CMS $\sigma_{tt}^{\rm tot}$        &   0.21      &   0.67  &    0.20      &   0.59  \\
CMS $t\bar{t}$ rap       &   0.96      &   0.96
&    0.94      &   0.96 \\
\hline 
LHCb total      &   {\bf 1.48}      &  {\bf 1.77 }
&  {\bf  1.47}      &   {\bf 1.62}  \\
LHCb $Z$ 940 pb      &   1.31     &   1.08 &    1.49      &   1.27  \\
LHCb $Z\to ee$ 2 fb      &   1.47     &   1.66
 &    1.14      &   1.33 \\
LHCb $W,Z \to \mu$ 7 TeV       &   1.54      &   1.51
&    1.76      &   1.60 \\
LHCb $W,Z \to \mu$ 8 TeV       &   1.51     &   2.28
 &    1.37     &   1.88 \\
\hline
\hline
Total dataset \bf     &\bf   1.187    & \bf  1.197   &  \bf  1.148      & \bf  1.168 \\
\hline
\end{tabular}
\caption{\small
\label{tab:chi2tab_31-nlo-nnlo-30-pc} 
Same as Tab.~\ref{tab:chi2tab_31-nlo-nnlo-30}, but now comparing the
default NNPDF3.1 NNLO and NNLO sets to the  variant in which charm is
perturbatively generated. For HERA $\sigma_c^{\rm NC}$ the number in
parenthesis refer to the subset of data to which the NNLO FC cut of
Table.~\ref{eq:tablekincuts} is applied.
}
\end{center}
\end{table}
%%%%%%%%%%%%%%%%%%%%%%%%%%%

%%%%%%%%%%%%%%%%%%%%%

In Table~\ref{tab:chi2tab_31-nlo-nnlo-30-pc} we show the 
$\chi^2/N_{\rm dat}$ values when charm is  perturbatively
generated  at NLO and NNLO. Unsurprisingly the fit quality deteriorates when
charm is not independently parametrized PDF. 
This is what one would naively expect 
since perturbative charm imposes a constraint upon the fit, thereby reducing the number
of free parameters. 

 However, it is interesting to observe that 
the fit quality to the inclusive HERA data (1306 data points)
significantly deteriorates
 when going from NLO to NNLO with perturbative charm, whereas it
 remains stable when charm is independently parametrized. Concerning
 the charm structure function data, note that,
 as discussed in Sect.~\ref{sec:datadis} above, a further cut is applied to 
 the HERA $\sigma_c^{\rm NC}$ data at NNLO when charm is independently
 parametrized. In order to allow for a consistent comparison, in
 Table~\ref{tab:chi2tab_31-nlo-nnlo-30-pc} we 
 show in parenthesis the value of 
 $\chi^2/N_{\rm dat}$ computed for the 37 (out of 47) data points that survive
 this cut also for all other  cases. Hence, for this data the fit
 quality is simlar with perturbative and parametrized charm, and also
 similar at NLO and NNLO (slightly worse at NNLO, by an amount
 compatible with a statistical fluctuation).
The fact  that when  parametrizing charm 
 there no longer is a deterioration of fit quality when going
from NLO to NNLO suggests that this  
resolves a tension present at NNLO, with  perturbative charm, 
between HERA and hadron collider data. Likewise, a purely perturbative charm leads
to a substantial deterioration at NNLO for BCDMS, NMC and especially
for the NuTeV dimuon cross-sections. This can be traced to the fact that
independently parametrizing charm is essential to reconcile the HERA data with the
constraints on the strange content of the proton imposed by
the ATLAS $W,Z$ 2011 rapidity distributions. 

In Fig.~\ref{fig:31-nnlo-fitted-vs-pch} we directly compare the PDFs
with parametrized and perturbative charm. The light quark PDFs and the
gluon are generally enhanced for $x\gsim0.003$ and reduced for smaller $x$
when charm is independently parametrized. The largest differences can be seen in the up quark, while
the strange and gluon distributions are more stable. The best-fit charm distribution has a
distinctly different shape and significantly larger uncertainty than
its perturbatively generated counterpart. As argued in
Ref.~\cite{Ball:2016neh} this shape might well be compatible with 
a charm PDF generated perturbatively at high perturbative orders. 

%%%%%%%%%%%%%%%%%%%%%%%%%%%%%%%%%%%%%%%%%%%%%%%%%%%%%%%%%%%%%%%%%%%%%
\begin{figure}[t]
  \begin{center}
    \includegraphics[scale=0.32]{images/plots/xu-31-nnlo-fitted-vs-pch.pdf}
    \includegraphics[scale=0.32]{images/plots/xd-31-nnlo-fitted-vs-pch.pdf}
    \includegraphics[scale=0.32]{images/plots/xubar-31-nnlo-fitted-vs-pch.pdf}
    \includegraphics[scale=0.32]{images/plots/xdbar-31-nnlo-fitted-vs-pch.pdf}
    \includegraphics[scale=0.32]{images/plots/xs-31-nnlo-fitted-vs-pch.pdf}
    \includegraphics[scale=0.32]{images/plots/xsbar-31-nnlo-fitted-vs-pch.pdf}
    \includegraphics[scale=0.32]{images/plots/xc-31-nnlo-fitted-vs-pch.pdf}
    \includegraphics[scale=0.32]{images/plots/xg-31-nnlo-fitted-vs-pch.pdf}
  \caption{\small 
    Comparison of NNPDF3.1 NNLO PDFs to a variant 
in which charm is generated entirely perturbatively 
(and everything else is unchanged).
    \label{fig:31-nnlo-fitted-vs-pch}
  }
\end{center}
\end{figure}
%%%%%%%%%%%%%%%%%%%%%%%%%%%%%%%%%%%%%%%%%%%%%%%%%%%%%%%%%%%%%%%%%%%%%%

In Fig.~\ref{fig:31-nnlo-fitted-vs-pch-err} we directly compare
PDF uncertainties. It is remarkable that the uncertainties other than
for charm are essentially unchanged when charm is independently parametrized,
with only a slight increase in sea quark PDF uncertainties for
$10^{-3}\lsim x \lsim 10^{-2}$. The uncertainty on the gluon is 
almost completely unaffected. The PDF uncertainty on charm when it is
independently parametrized  is in
line with that of other sea quark PDFs, while the uncertainty of the
perturbatively generated charm follows that of the gluon and is consequently 
much smaller.

%%%%%%%%%%%%%%%%%%%%%%%%%%%%%%%%%%%%%%%%%%%%%%%%%%%%%%%%%%%%%%%%%%%%%
\begin{figure}[t]
  \begin{center}
    \includegraphics[scale=0.32]{images/plots/xu-ERR-31-nnlo-fitted-vs-pch.pdf}
    \includegraphics[scale=0.32]{images/plots/xd-ERR-31-nnlo-fitted-vs-pch.pdf}
    \includegraphics[scale=0.32]{images/plots/xubar-ERR-31-nnlo-fitted-vs-pch.pdf}
    \includegraphics[scale=0.32]{images/plots/xdbar-ERR-31-nnlo-fitted-vs-pch.pdf}
    \includegraphics[scale=0.32]{images/plots/xs-ERR-31-nnlo-fitted-vs-pch.pdf}
    \includegraphics[scale=0.32]{images/plots/xsbar-ERR-31-nnlo-fitted-vs-pch.pdf}
    \includegraphics[scale=0.32]{images/plots/xc-ERR-31-nnlo-fitted-vs-pch.pdf}
    \includegraphics[scale=0.32]{images/plots/xg-ERR-31-nnlo-fitted-vs-pch.pdf}
  \caption{\small 
    Comparison of the fractional
    one-sigma PDF uncertainties in
    NNPDF3.1 NNLO with the corresponding version 
where charm is generated perturbatively
    (and everything else is unchanged).
    The PDF comparison plot was shown in
    Fig.~\ref{fig:31-nnlo-fitted-vs-pch}.
    \label{fig:31-nnlo-fitted-vs-pch-err}
  }
\end{center}
\end{figure}
%%%%%%%%%%%%%%%%%%%%%%%%%%%%%%%%%%%%%%%%%%%%%%%%%%%%%%%%%%%%%%%%%%%%%%

A previous comparison of PDFs determined with parametrized or perturbative
charm was presented in Ref.~\cite{Ball:2016neh} and led to the
conclusion that parametrizing charm and determining it from the data
 greatly reduces the dependence on the
charm mass thereby reducing the overall PDF uncertainty when the
uncertainty due to the charm mass is kept into account. 
As mentioned, NNPDF3.1 PDFs are determined using heavy quark
pole mass values and uncertainties recommended by the
Higgs Cross-Section Working Group~\cite{deFlorian:2016spz}.
For charm, this corresponds to
$ m_c^{\rm pole}=1.51 \pm 0.13\,{\rm GeV}$.
%
In order to estimate the impact of this uncertainty, we have produced 
NNPDF3.1 NNLO sets with $m_c^{\rm pole}=1.38$~GeV and $m_c^{\rm pole}=1.64$~GeV.
Results are shown in Fig.~\ref{fig:nnlo-fc-mcdep-1} for some
representative PDFs, both for the default NNPDF3.1 and for the version
with perturbative charm. It is clear that the very strong
dependence of the charm PDF on $m_c$ which is found when
charm is perturbatively generated all but disappears when charm is
independently parametrized. While the gluon is always quite stable, the 
dependence of perturbatively generated charm on $m_c$ propagates to
the light quark distributions. These are therefore significantly
stabilized by parametrizing charm. Indeed, if charm is generated
perturbatively, the shift in up and down quark distribution upon one-sigma
variation of the charm mass is comparable to (though somewhat smaller
than) 
the PDF uncertainty.
When charm is independently parametrized 
this dependence is considerably reduced. 
With parametrized charm, collider
observables at high scales become essentially independent of the charm mass, 
in line with the  expectation from decoupling arguments. 




%%%%%%%%%%%%%%%%%%%%%%%%%%%%%%%%%%%%%%%%%%%%%%%%%%%%%%%%%%%%%%%%%%%%%
\begin{figure}[t]
\begin{center}
  \includegraphics[scale=0.36]{images/plots/xc-31-nnlo-fc-mcdep.pdf}
  \includegraphics[scale=0.36]{images/plots/xc-31-nnlo-pc-mcdep.pdf}
  \includegraphics[scale=0.36]{images/plots/xg-31-nnlo-fc-mcdep.pdf}
  \includegraphics[scale=0.36]{images/plots/xg-31-nnlo-pc-mcdep.pdf}
  \includegraphics[scale=0.36]{images/plots/xu-31-nnlo-fc-mcdep.pdf}
  \includegraphics[scale=0.36]{images/plots/xu-31-nnlo-pc-mcdep.pdf}
  \includegraphics[scale=0.36]{images/plots/xd-31-nnlo-fc-mcdep.pdf}
  \includegraphics[scale=0.36]{images/plots/xd-31-nnlo-pc-mcdep.pdf}
  \caption{\small Dependence of the NNPDF3.1 NNLO PDFs on the charm
    mass. Results are shown both for parametrized charm  (left)
    and  perturbative charm (right), for  (from top to bottom) charm,
    gluon, up and down PDFs.
        \label{fig:nnlo-fc-mcdep-1}
  }
\end{center}
\end{figure}
%%%%%%%%%%%%%%%%%%%%%%%%%%%%%%%%%%%%%%%%%%%%%%%%%%%%%%%%%%%%%%%%%%%%%%








\clearpage

% Discussion on theoretical uncertainties
\subsection{Theoretical uncertainties}
\label{sec:thunc}

PDF uncertainties on global PDF sets entering the PDF4LHC15
combination
 consist only of the uncertainty propagated
from experimental data and uncertainties due to the
methodology. These
can be controlled through closure testing. There are however further sources
of uncertainty due to the theory used in PDF determination, which we
briefly assess here. These can be divided into two main classes:
\begin{itemize}
\item  Missing higher order uncertainties (MHOU), arising due to the truncation of
  the QCD perturbative expansion at a given fixed order (LO, NLO or NNLO)
  in the theory used for PDF determination.
\item Parametric uncertainties, due to the uncertainties on the values
  of parameters of the theory used for PDF determination: the main ones
 are the values of $\alpha_s(m_Z)$ and
  of $m_c^{\rm pole}$.
\end{itemize}

A full assessment of MHOU is an open problem, which we leave to future
investigations. For the time being, a first assessment can be obtained by
studying the perturbative stability of our results.
In Fig.~\ref{fig:distances_perturbativeorder} we show the distances at $Q=100$ GeV 
between all the PDFs in the
    LO and NLO sets, and in the NLO and NNLO sets. Some of the
 LO, NLO and NNLO PDFs are then compared directly in
 Fig.~\ref{fig:pdfcomp-pertorder}.  
Differences between the LO and NLO sets are very large, both for
central values and uncertainties, the latter being substantial at LO
due to the poor 
    fit quality. The shift in quark PDFs can be as large as two sigma
    ($d\simeq 20$), while the gluon at small $x$ is completely
    different between LO and NLO due to the fact that the singular small-$x$
    behaviour of the quark to gluon splittings only starts  at NLO, 
and due to the vanishing of gluon initiated  
    DIS and DY processes at LO.
On the other hand, when going from NLO to NNLO,
 PDF uncertainties are essentially unaffected. Central values are also 
 reasonably stable: the largest shifts, in the large-$x$ gluon and down quark and
 small-$x$ gluon, remain at or below the one-sigma level.

%%%%%%%%%%%%%%%%%%%%%%%%%%%%%%%%%%%%%%%%%%%%%%%%%%%%%%%%%%%%%%%%%%%%%
\begin{figure}[t]
\begin{center}
  \includegraphics[scale=1]{images/plots/distances_lo_vs_nlo.pdf}
  \includegraphics[scale=1]{images/plots/distances_nlo_vs_nnlo.pdf}
  \caption{\small Distances between the
    LO and NLO (top) and the NLO and NNLO (bottom) NNPDF3.1
    NNLO PDFs
    at $Q=100$ GeV. Note the difference in scale on the $y$ axis between
    the two plots.
    \label{fig:distances_perturbativeorder}
  }
\end{center}
\end{figure}
%%%%%%%%%%%%%%%%%%%%%%%%%%%%%%%%%%%%%%%%%%%%%%%%%%%%%%%%%%%%%%%%%%%%%%


%%%%%%%%%%%%%%%%%%%%%%%%%%%%%%%%%%%%%%%%%%%%%%%%%%%%%%%%%%%%%%%%%%%%%
\begin{figure}[t]
\begin{center}
  \includegraphics[scale=0.33]{images/plots/xg-pertorder.pdf}
  \includegraphics[scale=0.33]{images/plots/xu-pertorder.pdf}
  \includegraphics[scale=0.33]{images/plots/xdbar-pertorder.pdf}
  \includegraphics[scale=0.33]{images/plots/xsp-pertorder.pdf}
  \caption{\small Comparison between some 
    of the LO, NLO and NNPDF3.1 NNLO PDFs: gluon and up (top),
    antidown and total strangeness (bottom). All results are shown at
    $Q=100$ GeV, normalized to the NNLO central value.}
    \label{fig:pdfcomp-pertorder}
\end{center}
\end{figure}
%%%%%%%%%%%%%%%%%%%%%%%%%%%%%%%%%%%%%%%%%%%%%%%%%%%%%%%%%%%%%%%%%%%%%%

A quantitative estimate of the MHOU can be obtained
by computing the shift between the central values of the NLO and NNLO
NNPDF3.1 PDFs. The result is shown in Fig.~\ref{fig:THerror-shifts}
for some PDF combinations. In the plot, the shift has been
symmetrized, and is compared to the NLO standard PDF uncertainty.
In the quark singlet $\Sigma$ for $x\lesssim 10^{-3}$ the shift is
larger than the PDF uncertainty, while it is smaller for 
individual flavors (as illustrated by the two quark
distributions shown).  
This suggests that for individual quark flavors and the gluon at
NNLO, MHOU can be reasonably neglected at the current level of precision. 
However, for particular combinations (such as the singlet at small $x$) it is
unclear whether MHOU can be neglected even at NNLO, given that at NLO
they are larger than the PDF uncertainty.

%
%%%%%%%%%%%%%%%%%%%%%%%%%%%%%%%%%%%%%%%%%%%%%%%%%%%%%%%%%%%%%%%%%%%%%
\begin{figure}[t]
\begin{center}
  \includegraphics[scale=0.36]{images/plots/xsinglet-THerror.pdf}
  \includegraphics[scale=0.36]{images/plots/xg-THerror.pdf}
  \includegraphics[scale=0.36]{images/plots/xu-THerror.pdf}
  \includegraphics[scale=0.36]{images/plots/xdbar-THerror.pdf}
  \caption{\small Comparison between the NLO PDF uncertainties and the
    shift between the NLO and NNLO PDFs. All results are shown as
    ratios to the NLO PDFs, for $Q=100$~GeV. The shift is
    symmetrized. We show results for the singlet, gluon (top); up and
    antidown (bottom) PDFs.
    \label{fig:THerror-shifts}
  }
\end{center}
\end{figure}
%%%%%%%%%%%%%%%%%%%%%%%%%%%%%%%%%%%%%%%%%%%%%%%%%%%%%%%%%%%%%%%%%%%%%%

We finally turn to parametric uncertainties. As we have discussed in
Sect.~\ref{sec:results-mc}, the dependence of PDFs upon the charm mass 
is almost entirely removed by parametrizing charm. The dependence on the $b$-quark 
mass is minor, except for the bottom PDFs
themselves~\cite{Ball:2011mu,Harland-Lang:2015qea}. Therefore, the only
significant residual parametric uncertainty is on the value of the strong
coupling. This uncertainty is routinely included along with the PDF
uncertainty; in order to do this consistently, one needs PDF sets
produced with different central values of $\alpha_s$ (see
e.g. Ref.~\cite{Butterworth:2015oua}).  We have 
determined  NNPDF3.1 NLO and NNLO  PDFs with $\alpha_s(m_Z)$
varied in the range
$0.108\le\alpha_s(m_Z)\le0.124$ (see Sect.~\ref{sec:delivery}).

%%%%%%%%%%%%%%%%%%%%%%%%%%%%%%%%%%%%%%%%%
\begin{figure}[t]
  \begin{center}
  \includegraphics[scale=0.36]{images/plots/xg-31-nlo-fc-asdep.pdf}
  \includegraphics[scale=0.36]{images/plots/xu-31-nlo-fc-asdep.pdf}
  \includegraphics[scale=0.36]{images/plots/xg-31-nnlo-fc-asdep.pdf}
  \includegraphics[scale=0.36]{images/plots/xu-31-nnlo-fc-asdep.pdf}
  \caption{\small 
    \label{fig:PDFasvariations} Dependence of NNPDF3.1 NLO (top) and
    NNLO (bottom) PDFs on
    the value of $\alpha_s$. The gluon (left) and up quark (right) are
    shown at $Q=100$~GeV, normalized to the central value.
  }
\end{center}
\end{figure}
%%%%%%%%%%%%%%%%%%%%%%%%%%%%%%%%%%%%%%%%

In Fig.~\ref{fig:PDFasvariations} we compare the up and gluon PDFs as
$\alpha_s(m_Z)$ is varied by $\Delta\alpha_s=\pm0.002$ about its
central value. As is well known,
the gluon is anti-correlated to $\alpha_s(m_Z)$ at small and medium $x$, but 
positively correlated to it at large $x$. The dependence on $\alpha_s$ is 
rather milder for quark PDFs, with positive correlation at small $x$, and very little
dependence altogether at large $x$.

   



\clearpage
